% GENERATED FILE. Edit canonical lessons, then run npm run generate:books.
% canonical-source-hash: fnv1a64:f93ef0c6b8438f8b
% canonical-lessons: JA-C310-zutto, JA-C310-jitsu-wa, JA-C310-toutou, JA-C310-kanau, JA-C310-mamonaku

\chapter{Waiting for an Answer}
\label{ch:ja-waiting-for-an-answer}
\hlchaptermodality{310}

\begin{chapteropening}
\textbf{By the end of this chapter:} \emph{I can say how long I waited and what came true.}

Complete the last lesson of chapter 310: I can say how long I waited and what came true.
\end{chapteropening}

\section[\texorpdfstring{\ja{ずっと} (zutto)}{ずっと (zutto)}]{\ja{ずっと} (zutto) — all along; the whole time}

\label{lesson:JA-C310-zutto}



\begin{quote}
Before the new one: say the Japanese for a person, then the Japanese for to confide.
\end{quote}

\subsection*{You'll want to know: \ja{ずっと}}
\textbf{\ja{ずっと}} — \emph{zutto} — \textquotedblleft{}all along; the whole time\textquotedblright{}.

\begin{quote}
\textbf{\ja{ずっと} \ja{へんじを} \ja{まっています。}} — \emph{zutto henji o matte imasu} — \textquotedblleft{}I have been waiting for an answer all this time.\textquotedblright{}
\end{quote}

\begin{cousinweb}
\textbf{\ja{ずっと}} (\emph{zutto}) is a native Japanese word (wago), inherited rather than borrowed from Chinese, and it is usually written in kana (JMdict).
\end{cousinweb}

\begin{grammarlens}[title={What you've built}]
One more word for telling what happened.
\end{grammarlens}

\subsection*{Your turn}
\begin{itemize}
\raggedright
  \item \emph{Say it:} \emph{zutto}
  \item \emph{Say it:} \emph{zutto}, once more
  \item \emph{Say it:} \emph{zutto henji o matte imasu} — I have been waiting for an answer all this time.
  \item \emph{Recall:} say the Japanese for a person, then the Japanese for to confide, then say \emph{zutto} again
\end{itemize}

\begin{tcolorbox}[breakable,colback=teal!4,colframe=teal!35!black,title={Before you move on}]
What does \ja{ずっと} mean? (\textquotedblleft{}All along; the whole time\textquotedblright{}.) Say it once more.
\end{tcolorbox}
\section[\texorpdfstring{\ja{じつは} (jitsu wa)}{じつは (jitsu wa)}]{\ja{じつは} (jitsu wa) — actually; to tell the truth}

\label{lesson:JA-C310-jitsu-wa}



\begin{quote}
Before the new one: say the Japanese for to confide, then the Japanese for all along.
\end{quote}

\subsection*{You'll want to know: \ja{じつは}}
\textbf{\ja{じつは}} — \emph{jitsu wa} — \textquotedblleft{}actually; to tell the truth\textquotedblright{}.

\begin{quote}
\textbf{\ja{じつは、きのう} \ja{びょういんに} \ja{いきました。}} — \emph{jitsu wa, kinō byōin ni ikimashita} — \textquotedblleft{}To tell the truth, I went to the hospital yesterday.\textquotedblright{}
\end{quote}

\begin{cousinweb}
\textbf{\ja{じつは}} (\emph{jitsu wa}) is a mixed word, native and Sino-Japanese: \emph{jitsu}, \textquotedblleft{}truth\textquotedblright{}, is a Chinese character read with its on'yomi, a reading borrowed from Chinese (Middle Chinese \emph{zyit}), and \emph{wa}, \textquotedblleft{}the topic particle\textquotedblright{}, is native Japanese (JMdict, KANJIDIC2; Middle Chinese in Baxter's transcription).
\end{cousinweb}

\begin{grammarlens}[title={What you've built}]
One more word for telling what happened.
\end{grammarlens}

\subsection*{Your turn}
\begin{itemize}
\raggedright
  \item \emph{Say it:} \emph{jitsu wa}
  \item \emph{Say it:} \emph{jitsu wa}, once more
  \item \emph{Say it:} \emph{jitsu wa, kinō byōin ni ikimashita} — To tell the truth, I went to the hospital yesterday.
  \item \emph{Recall:} say the Japanese for to confide, then the Japanese for all along, then say \emph{jitsu wa} again
\end{itemize}

\begin{tcolorbox}[breakable,colback=teal!4,colframe=teal!35!black,title={Before you move on}]
What does \ja{じつは} mean? (\textquotedblleft{}Actually; to tell the truth\textquotedblright{}.) Say it once more.
\end{tcolorbox}
\section[\texorpdfstring{\ja{とうとう} (tōtō)}{とうとう (tōtō)}]{\ja{とうとう} (tōtō) — in the end; at long last}

\label{lesson:JA-C310-toutou}



\begin{quote}
Before the new one: say the Japanese for all along, then the Japanese for actually.
\end{quote}

\subsection*{You'll want to know: \ja{とうとう}}
\textbf{\ja{とうとう}} — \emph{tōtō} — \textquotedblleft{}in the end; at long last\textquotedblright{}.

\begin{quote}
\textbf{\ja{とうとう} \ja{へんじは} \ja{きませんでした。}} — \emph{tōtō henji wa kimasen deshita} — \textquotedblleft{}In the end no answer came.\textquotedblright{}
\end{quote}

\begin{cousinweb}
\textbf{\ja{とうとう}} (\emph{tōtō}) is usually written in kana (JMdict). JMdict also gives a spelling it marks as rare: two Chinese characters read with their on'yomi, readings borrowed from Chinese, \emph{tō}, \textquotedblleft{}arrival\textquotedblright{} (Middle Chinese \emph{tawH}) and \emph{tō}, \textquotedblleft{}head\textquotedblright{} (\emph{duw}) (KANJIDIC2; Middle Chinese in Baxter's transcription). Whether the word is Sino-Japanese, built from them, or a native adverb later given that spelling is not settled.
\end{cousinweb}

\begin{grammarlens}[title={What you've built}]
One more word for telling what happened.
\end{grammarlens}

\subsection*{Your turn}
\begin{itemize}
\raggedright
  \item \emph{Say it:} \emph{tōtō}
  \item \emph{Say it:} \emph{tōtō}, once more
  \item \emph{Say it:} \emph{tōtō henji wa kimasen deshita} — In the end no answer came.
  \item \emph{Recall:} say the Japanese for all along, then the Japanese for actually, then say \emph{tōtō} again
\end{itemize}

\begin{tcolorbox}[breakable,colback=teal!4,colframe=teal!35!black,title={Before you move on}]
What does \ja{とうとう} mean? (\textquotedblleft{}In the end; at long last\textquotedblright{}.) Say it once more.
\end{tcolorbox}
\section[\texorpdfstring{\ja{かなう} (kanau)}{かなう (kanau)}]{\ja{かなう} (kanau) — to come true}

\label{lesson:JA-C310-kanau}



\begin{quote}
Before the new one: say the Japanese for actually, then the Japanese for in the end.
\end{quote}

\subsection*{You'll want to know: \ja{かなう}}
\textbf{\ja{かなう}} — \emph{kanau} — \textquotedblleft{}to come true\textquotedblright{}.

\begin{quote}
\textbf{\ja{いつか} \ja{ゆめが} \ja{かなうと} \ja{おもいます。}} — \emph{itsuka yume ga kanau to omoimasu} — \textquotedblleft{}I think my dream will come true one day.\textquotedblright{}
\end{quote}

Polite \textbf{\ja{かないます}}, \emph{kanaimasu}; a godan verb. Past (polite): \textbf{\ja{ゆめが} \ja{かないました。}}, \emph{yume ga kanaimashita}.

\begin{cousinweb}
\textbf{\ja{かなう}} (\emph{kanau}) is a native Japanese word (wago), inherited rather than borrowed from Chinese; its kanji was attached later as a kun'yomi (JMdict).
\end{cousinweb}

\begin{grammarlens}[title={What you've built}]
One more word for telling what happened.
\end{grammarlens}

\subsection*{Your turn}
\begin{itemize}
\raggedright
  \item \emph{Say it:} \emph{kanau}
  \item \emph{Say it:} \emph{kanau}, once more
  \item \emph{Say it:} \emph{itsuka yume ga kanau to omoimasu} — I think my dream will come true one day.
  \item \emph{Recall:} say the Japanese for actually, then the Japanese for in the end, then say \emph{kanau} again
\end{itemize}

\begin{tcolorbox}[breakable,colback=teal!4,colframe=teal!35!black,title={Before you move on}]
What does \ja{かなう} mean? (\textquotedblleft{}To come true\textquotedblright{}.) Say it once more.
\end{tcolorbox}
\section[\texorpdfstring{\ja{まもなく} (mamonaku)}{まもなく (mamonaku)}]{\ja{まもなく} (mamonaku) — shortly; in a moment}

\label{lesson:JA-C310-mamonaku}



\begin{quote}
Before the new one: say the Japanese for in the end, then the Japanese for to come true.
\end{quote}

\subsection*{You'll want to know: \ja{まもなく}}
\textbf{\ja{まもなく}} — \emph{mamonaku} — \textquotedblleft{}shortly; in a moment\textquotedblright{}.

\begin{quote}
\textbf{\ja{まもなく} \ja{でんしゃが} \ja{きます。}} — \emph{mamonaku densha ga kimasu} — \textquotedblleft{}The train will be here shortly.\textquotedblright{}
\end{quote}

\begin{cousinweb}
\textbf{\ja{まもなく}} (\emph{mamonaku}) is a native Japanese word (wago), inherited rather than borrowed from Chinese, built from \emph{ma}, \textquotedblleft{}an interval\textquotedblright{}, \emph{mo}, \textquotedblleft{}even\textquotedblright{}, and \emph{naku}, \textquotedblleft{}without\textquotedblright{} (from \emph{nai}) (JMdict, KANJIDIC2).
\end{cousinweb}

\begin{grammarlens}[title={What you've built}]
One more word for telling what happened.
\end{grammarlens}

\subsection*{Your turn}
\begin{itemize}
\raggedright
  \item \emph{Say it:} \emph{mamonaku}
  \item \emph{Say it:} \emph{mamonaku}, once more
  \item \emph{Say it:} \emph{mamonaku densha ga kimasu} — The train will be here shortly.
  \item \emph{Recall:} say the Japanese for in the end, then the Japanese for to come true, then say \emph{mamonaku} again
\end{itemize}

\begin{tcolorbox}[breakable,colback=teal!4,colframe=teal!35!black,title={Before you move on}]
What does \ja{まもなく} mean? (\textquotedblleft{}Shortly; in a moment\textquotedblright{}.) Say it once more.
\end{tcolorbox}
